\documentclass[11pt,a4paper]{article}

\usepackage[T1]{fontenc}
\usepackage[utf8]{inputenc}
\usepackage[french]{babel}
\usepackage{geometry}
\geometry{top=2.2cm,bottom=2.2cm,left=1.8cm,right=1.8cm}
\usepackage{lmodern}
\usepackage{amsmath,amssymb,amsthm}
\usepackage{booktabs}
\usepackage{tabularx}
\usepackage{enumitem}
\setlength{\parskip}{4pt}
\usepackage{titlesec}
\usepackage{fancyhdr}
\usepackage{hyperref}
\hypersetup{colorlinks=false,pdfborder={0 0 0}}
\usepackage{array}
\usepackage{microtype}
\usepackage{tikz}
\usetikzlibrary{arrows.meta,positioning,shapes.geometric}
\usepackage{pifont}
\usepackage{makecell}
\usepackage{caption}
\usepackage{float}
\setlength{\emergencystretch}{2.5em}
% ---------- typography ----------
\titleformat{\section}
  {\normalsize\bfseries\uppercase}
  {\thesection.}{0.5em}{}[\vspace{-2pt}\rule{\linewidth}{0.4pt}\vspace{2pt}]
\titleformat{\subsection}
  {\normalsize\bfseries}{\thesubsection.}{0.5em}{}
\titlespacing{\section}{0pt}{10pt}{4pt}
\titlespacing{\subsection}{0pt}{7pt}{2pt}

% ---------- header ----------
\pagestyle{fancy}
\fancyhf{}
\renewcommand{\headrulewidth}{0.4pt}
\fancyhead[L]{\small\textit{Research Statement --- CIFRE PhD}}
\fancyhead[R]{\small\textit{Yassine Mannai}}
\fancyfoot[C]{\small\thepage}

% ---------- theorem-like envs ----------
\newtheorem{rqthm}{Research Question}

\begin{document}

% ================================================================
%  TITLE BLOCK
% ================================================================
\begin{center}
{\Large\bfseries Research Statement}\\[4pt]
{\large\bfseries CIFRE PhD --- CREST - IMT, Generali France}\\[8pt]
{\normalsize\itshape
Regime-Conditional Factor Dynamics, Functional Capital Proxies\\
and Optimal Stochastic Control for Insurance Portfolio Management}
\end{center}

\medskip
\noindent
\begin{tabular}{@{}l>{\raggedright\arraybackslash}p{11.8cm}@{}}
\textbf{Candidate:}            & Yassine Mannai \\
\textbf{Contact:}              & \texttt{yassine.mannai@dauphine.eu}
                                 $\;\cdot\;$ +33\,(0)6\,23\,64\,31\,80 \\
\textbf{Programme:}            & MSc in Quantitative Finance (272),
                                 Paris Dauphine-PSL \\
\textbf{Academic supervisors:} & Caroline Hillairet (CREST, ENSAE Paris) and \\
                               & Anthony R\'{e}veillac (INSA Toulouse, IMT) \\
\textbf{Industry supervisor:}  & Mamy Ramamonjisoa ---
                                 SAA team, Generali France \\
\textbf{Investment Division:}  & Cedrik d'Aviau de Ternay --- Chief Investment Officer, Generali France \\
\textbf{Start date:}           & January 2027
                                 $\;\cdot\;$ Duration: 36 months
\end{tabular}

\noindent\rule{\linewidth}{0.8pt}

\medskip
\noindent\textit{\textbf{Thesis in one sentence.}}
\begin{quote}\itshape
This thesis introduces the first \emph{regime-conditional generative framework
with certified $W_2$ guarantees under online regime inference}, coupled to a
\emph{differentiable functional capital proxy}, to enable a fully dynamic,
capital-aware Total Portfolio Approach for a life insurer under Solvency~II.
The three axes form a single closed control loop: a certified scenario
generator (Axis~1) feeds a fast differentiable SCR proxy (Axis~2), which
becomes the capital constraint of a joint allocation-and-hedging policy
(Axis~3).
\end{quote}

% ================================================================
\section{Industrial Context and Motivation}
% ================================================================

\subsection{The 2022 dependence crisis and the limits of static allocation}

The year 2022 was a structural stress test for insurance asset managers: unlike
earlier single-asset shocks, equities and bonds fell simultaneously and
diversification broke down. For a life insurer such as Generali France, this
exposed a structural flaw in Strategic Asset Allocation (SAA), which assumes
stable, linear dependence between economic factors. The equity/bond correlation
switched from $-0.3$ (diversifying) to $+0.6$ (concentration risk) within
weeks --- a regime reversal that fixed-correlation models cannot anticipate.

This motivates the central problem of this thesis:
\textit{how to characterise and optimally control the solvency position of an insurer
when the joint dynamics of economic factors are governed by a latent Markov chain
inducing non-stationary, non-linear dependence structures across regimes?}

\subsection{Three compounded bottlenecks in the current SAA process}

The SAA process at Generali France runs in three sequential steps, each with
structural limitations.
\textbf{(1)~Economic Scenario Generator (ESG)} --- a proprietary ESG produces
thousands of trajectories (rates, spreads, equity, real estate, inflation)
calibrated with \textbf{fixed correlations averaged across regimes}, feeding all
downstream actuarial calculations.
\textbf{(2)~SCR via nested LSMC} --- the Solvency Capital Requirement, the
$99.5\%$ one-year VaR of net asset value change, requires the Best Estimate of
Liabilities (BEL) under $10{,}000+$ scenarios; Least-Squares Monte Carlo
(LSMC, Krah et al., 2018) replaces the inner simulation by an offline
polynomial proxy.
\textbf{(3)~Annual Markowitz optimisation} --- the SCR is a fixed constraint in
a mean-variance optimisation revised \textbf{once per year}, with no automatic
identification of the dominant portfolio objective (duration, inflation, carry,
capital efficiency), even though the relevant factor varies across portfolios
and regimes.

\medskip
\noindent\textbf{Four structural limitations.}

\begin{enumerate}[leftmargin=1.8em,itemsep=2pt]
\item \textbf{Regime-blind dependence.}
Fixed correlations cannot anticipate regime reversals: a Markowitz optimum for
a low-inflation regime becomes high-risk under stagflation.

\item \textbf{LSMC proxy limitations.}
The polynomial proxy is accurate only within its regime-agnostic calibration
region, degrading precisely in the stress zones most relevant for capital
decisions, and is not expressed in common economic factors.

\item \textbf{Static, siloed decision-making.}
Allocation and hedging are decided separately --- suboptimal for an insurer with
convex critical losses --- and the objective is imposed uniformly rather than
inferred from each portfolio's dominant risk factor.

\item \textbf{Private asset valuation bias.}
Private equity, private debt and real estate NAVs are smoothed quarterly: in
2008 private equity NAVs stayed stable for $6$--$12$ months before correcting
sharply, understating stress-regime dependence and introducing a lag bias in
the SCR.
\end{enumerate}

\subsection{The Total Portfolio Approach: target framework}

The \textbf{Total Portfolio Approach (TPA)} evaluates each decision through its
marginal contribution to total balance-sheet risk via common factors
(Duration, Inflation, Growth, Liquidity, Carry). Implementations at sovereign
wealth funds (Future Fund, CPP Investments) report about $+1\%$ excess return
per year over classical SAA (Thinking Ahead Institute, 2025; CAIA, 2024). For
an insurer under Solvency~II, the factor space must also include liability-side
factors (surrender, longevity, yield-curve/BEL interactions), and every
decision must internalise a capital cost. Existing actuarial tools address none
of this jointly; this thesis fills the gap along three coupled axes:
regime-conditional factor simulation (Axis~1), fast capital proxy (Axis~2),
and joint allocation-hedging policy (Axis~3).

% ================================================================
\section{Research Overview}
% ================================================================

\begin{figure}[H]
\centering
\begin{tikzpicture}[
  node distance=6mm,
  box/.style={draw,rounded corners,align=center,inner sep=4pt,
              minimum height=12mm,font=\small},
  ax/.style={box,minimum width=37mm},
  arr/.style={-{Latex[length=2mm]},thick},
  lbl/.style={font=\scriptsize\itshape,align=center}
]
\node[ax] (a1) {\textbf{Axis~1}\\Regime-conditional\\generative model\\\scriptsize(HMM-SMC + GAN/VAE/SGM)};
\node[ax,right=17mm of a1] (a2) {\textbf{Axis~2}\\Differentiable\\SCR proxy\\\scriptsize(Volterra functional)};
\node[ax,right=17mm of a2] (a3) {\textbf{Axis~3}\\Joint allocation\\\& hedging policy\\\scriptsize(stochastic control)};
\draw[arr] (a1) -- node[lbl,above]{certified\\scenarios ($W_2$)} (a2);
\draw[arr] (a2) -- node[lbl,above]{capital\\constraint ($\nabla$SCR)} (a3);
\draw[arr] (a3.south) -- ++(0,-7mm) -| node[lbl,below,pos=0.25]{realised regime / rebalancing feedback} (a1.south);
\end{tikzpicture}
\caption{The three axes form a single closed capital-aware control loop:
certified regime-conditional scenarios feed a fast differentiable capital
proxy, which becomes the explicit constraint of the joint allocation and
hedging policy; realised regimes feed back into the generator.}
\label{fig:pipeline}
\vspace{4pt}

\renewcommand{\arraystretch}{1.15}
\begin{tabularx}{\linewidth}{@{}Xcccc@{}}
\toprule
\textbf{Approach} & \makecell{Regime-\\aware} & \makecell{Differen-\\tiable} &
\makecell{$W_2$/$L^2$\\certified} & \makecell{Insurance-\\calibrated} \\
\midrule
Static SAA / Markowitz                       & \ding{55} & \ding{55} & \ding{55} & \ding{51} \\
Copulas (Sklar, Fermanian)                   & \ding{55} & \ding{55} & \ding{51} & \ding{55} \\
LSMC proxy (Krah et al.)                     & \ding{55} & \ding{55} & \ding{51} & \ding{51} \\
GAN/VAE time-series (Wiese, Yoon)            & \ding{55} & \ding{51} & \ding{55} & \ding{55} \\
Score-based models (Ocello et al.)           & \ding{55} & \ding{51} & \ding{51} & \ding{55} \\
Deep hedging (Buehler et al.)               & \ding{55} & \ding{51} & \ding{55} & \ding{55} \\
\midrule
\textbf{This thesis (Axes 1--3)}            & \ding{51} & \ding{51} & \ding{51} & \ding{51} \\
\bottomrule
\end{tabularx}
\captionof{table}{Positioning of the research programme against the state of the art.\\
\ding{51}~= requirement met, \ding{55}~= not addressed.}
\label{tab:positioning}
\end{figure}

% ================================================================
\section{Research Programme}
% ================================================================

% ----------------------------------------------------------------
\subsection{Axis~1 --- Regime-Conditional Factor Dependence}
% ----------------------------------------------------------------

\noindent\textbf{Research question.}
\begin{rqthm}
How to generate the joint distribution of TPA factors conditionally on a
sequentially estimated macroeconomic regime, while providing: (i)~formal
goodness-of-fit guarantees on the learned dependence structure; and
(ii)~non-asymptotic $W_2$-convergence bounds under particle-filter
inference error on the regime state?
\end{rqthm}

\noindent\textbf{State of the art.}

\textit{Regime identification.}
Hidden Markov Models (Hamilton, 1989) and Sequential Monte Carlo (SMC) particle
filters (Chopin \& Papaspiliopoulos, 2020) provide tractable real-time estimation
of latent macroeconomic regimes $z_t \in \{1,\ldots,K\}$. Their combination ---
an online HMM-SMC filter --- yields a posterior $p(z_t|\mathcal{F}_t)$ updated at
each observation.

\textit{Generative Adversarial Networks.}
Applied to financial series by Wiese et al.\ (2020, Quant-GAN) and Yoon et al.\
(2019, TimeGAN), conditional GANs generate
$\mathbf{F}_{t+1} | z_t = k \sim G_\theta(\varepsilon,k)$ with full flexibility
(no parametric dependence assumption), but suffer from \textit{mode collapse}
and the absence of distributional guarantees, limiting regulatory auditability.

\textit{Variational Autoencoders.}
CVAEs encode each trajectory into a regime-conditional posterior and maximise
the ELBO; they offer stable training and an auditable latent structure, but
produce overly smooth samples due to the Gaussian posterior approximation.

\textit{Score-Based Generative Models (SGMs).}
Strasman, Ocello et al.\ (TMLR, 2024) establish $W_2$ convergence bounds:
\begin{equation}
  \mathrm{KL}(\hat{p}\,\|\,p^*)
  \;\le\; \int_0^T g(t)^2\,\varepsilon_{\mathrm{score}}(t)\,dt
  + \varepsilon_{\mathrm{init}}.
\label{eq:kl}
\end{equation}
Gentiloni-Silveri \& Ocello (ICML, 2025) extend these guarantees to
non-log-concave distributions (Gaussian mixtures), directly applicable to
multi-regime market models. De Marco, Pham \& Zanni (2025) further introduce
jump-diffusion Schr\"{o}dinger bridges to capture discontinuous regime
transitions.

\textit{Copulas: a validation framework, not a generative model.}
Sklar's theorem separates margins from dependence, and Fermanian (2013, 2022)
provides goodness-of-fit tests for conditional copulas. Classical parametric
families (Gaussian, Student-$t$, Archimedean), however, impose rigid tails,
scale poorly in high dimension and cannot capture cross-regime dynamics.
Copulas therefore serve here exclusively as a \emph{validation tool}: the
conditional tail-dependence coefficient $\lambda_k$ is estimated and tested
via Fermanian's specification tests to validate generative outputs, not to
generate them.

\textit{Financial networks and correlation structure.}
Bardoscia et al.\ (2021) show that asset correlation matrices exhibit community
structure and spectral properties amenable to graph-theoretic analysis,
motivating an optional graph-neural-network (GNN) encoder in the Axis~1
generator to capture non-linear pairwise interactions beyond the covariance
matrix.

\noindent\textbf{Identified gaps.}
Neither conditional GANs, VAEs nor SGMs have been applied to the joint
distribution of TPA factors with real-time regime identification.
For the GAN/VAE: the absence of formal regime-conditional tail-dependence
validation. For the SGM: the absence of $W_2$ bounds for the conditional case
with sequentially estimated regime states. These constitute the theoretical
contribution of Axis~1.

\noindent\textbf{Proposed approach.}

\textit{Stage 1 (Year 1) --- cGAN and cVAE.}
Let $z_t$ be identified by an online HMM-SMC filter. A conditional GAN
$G_\theta(\varepsilon,k)$ and a CVAE are trained in parallel, both incorporating
the private asset vintage lag factor
$\tilde{F}^{\mathrm{PA}}_t = \mathrm{NAV}^{\mathrm{PA}}_{t-1}$
to capture the quarterly smoothing bias. Model comparison is conducted on:
Wasserstein-2 distance $W_2$, Maximum Mean Discrepancy (MMD), per-regime
tail calibration via Fermanian's $\lambda_k$ tests, autocorrelation structure,
and regulatory auditability.

\textit{Stage 2 (Year 2) --- Conditional SGM with $W_2$ guarantees.}
A conditional OU forward process for regime $k$:
\begin{equation}
  d\mathbf{x}^\tau = -\mathbf{x}^\tau\,d\tau
  + \sqrt{2}\,g_k(\tau)\,dW_\tau, \quad \tau \in [0,T],
\label{eq:sgm}
\end{equation}
with regime-specific noise schedule $g_k(\tau)$. The theoretical contribution
is the extension of the $W_2$ bounds of Gentiloni-Silveri \& Ocello (ICML 2025)
to the conditional case with online SMC regime estimation --- quantifying how
inference error in $p(z_t|\mathcal{F}_t)$ propagates into the generative bound.
A scenario selection layer then extracts a central regime-conditional scenario
for downstream reporting and optimisation diagnostics.

% ----------------------------------------------------------------
\subsection{Axis~2 --- Supervised Capital Proxy}
% ----------------------------------------------------------------

\noindent\textbf{Research question.}
\begin{rqthm}
How to construct a proxy for the insurance SCR --- expressed as a function of
TPA factors, the current regime, and liability states including private asset
NAV lags formalised as a Volterra functional --- that is accurate in
decision-relevant zones with controlled $L^2$ error bounds per regime and
sub-10ms inference time, and exposes a differentiable interface for its direct
embedding as an explicit capital constraint in a downstream stochastic
optimisation programme?
\end{rqthm}

\noindent\textbf{State of the art.}

\textit{Classical proxy models.}
The LSMC framework (Krah et al., 2018) approximates the BEL by polynomial
regression on a design-of-experiment. Its two critical limitations are
regime-blindness and degraded accuracy in stress zones. Neural network proxies
(Huber \& Schmocker, 2022) achieve $1$--$3\%$ accuracy at $<10$ms inference.
Bourgey (2020) provides theoretical foundations via multi-level Monte Carlo
(MLMC) estimators for nested expectations and polynomial chaos expansions with
$L^2$ error estimates, motivating the adaptive-weighting architecture of
Axis~2.

\textit{Recent generative and flow-based surrogates.}
Three recent architectures merit systematic comparison for their
uncertainty-quantification and stability properties:
\textbf{Neural ODEs} (Chen et al., 2018), which model the proxy surface as a
continuous-depth dynamical system; \textbf{Conditional Flow Matching} (Lipman
et al., 2022), which learns straight-line ODE paths offering stable training
and a full conditional density
$p(\widehat{\mathrm{SCR}} | \mathbf{F}, z, \mathbf{P})$ for uncertainty-aware
allocation; and \textbf{score-based diffusion surrogates}. \textbf{Tabular
diffusion models} (TabDDPM, Kotelnikov et al., 2023) additionally serve as
design-of-experiment augmentation for under-represented regimes.

\textit{Volterra processes and path-dependent risk measures.}
The private asset NAV lag is not a Markovian state variable. The observed
smoothed value
$\widetilde{\mathrm{NAV}}_t = \int_0^t K(t,s)\,\mathrm{NAV}^*_s\,ds$
is a Volterra functional of the true (unobservable) NAV process, where $K(t,s)$
is a smoothing kernel. The theory of affine Volterra processes (Abi Jaber,
Larsson \& Pulido, 2019; El Euch \& Rosenbaum, 2019) provides a tractable
framework for such path-dependent dynamics. This motivates extending the
capital risk measure to the functional setting:
\begin{equation}
  \rho^V(L_T)
  = \phi\!\left(\int_0^T K(T,s)\,\mathrm{CVaR}_\alpha(L_s)\,ds\right),
\label{eq:volterra}
\end{equation}
where $\phi$ is a monetary risk measure. Equation~\eqref{eq:volterra}
formalises the temporally diffused risk crystallisation of private assets and
provides a theoretical grounding for the vintage lag correction embedded in the
proxy state vector.

\textit{Model comparison and interoperability metrics.}
Proxy architectures will be evaluated on two sets of criteria.
\textit{Comparison metrics}: (i)~accuracy --- RMSE, relative error at the SCR
regulatory threshold, $L^\infty$ norm in stress zones; (ii)~calibration ---
prediction-interval coverage, reliability diagrams; (iii)~speed --- inference
latency (target $<5$ms) and batch throughput; (iv)~interpretability ---
SHAP values, integrated gradients, first-order Sobol sensitivity indices
$S_i = \mathrm{Var}[\mathbb{E}[\widehat{\mathrm{SCR}}|\mathbf{F}_i]]/
\mathrm{Var}[\widehat{\mathrm{SCR}}]$.
\textit{Interoperability metrics}: (i)~differentiability ---
$\nabla_\mathbf{F}\widehat{\mathrm{SCR}}$ available via automatic
differentiation for gradient-based embedding in the Axis~3 optimiser;
(ii)~API compatibility with GPS and PyTorch/JAX optimisation frameworks;
(iii)~monotonicity and convexity preservation as required by Solvency~II
constraints; (iv)~memory footprint for production deployment.

\noindent\textbf{Identified gaps.}
No existing proxy is jointly conditioned on the macroeconomic regime, TPA
factors and private asset NAV lag. None accounts for the path-dependent
Volterra structure of NAV smoothing, and none provides the differentiable
interface required for its direct embedding as an explicit capital constraint
in a dynamic optimisation programme.

\noindent\textbf{Proposed approach.}
We learn the supervised proxy:
\begin{equation}
  \widehat{\mathrm{SCR}}(t)
  = f_\theta\!\bigl(\mathbf{F}_t,\,z_t,\,\mathbf{P}_t\bigr),
\label{eq:proxy}
\end{equation}
where $\mathbf{P}_t$ includes effective duration, average guaranteed rate,
observed lapse ratio and the vintage lag $\tilde{F}^{\mathrm{PA}}_{t-1}$.
The network $f_\theta$ is trained with \textbf{adaptive sample weighting}
concentrating mass near the regulatory threshold and in stress regimes.
The generalisation bound:
\begin{equation}
  \bigl|\mathrm{SCR}(t) - f_\theta(\mathbf{F}_t,z_t,\mathbf{P}_t)\bigr|
  \;\le\; \varepsilon_k(n_k,\delta),
\label{eq:bound}
\end{equation}
decreasing with $n_k$ (training points in regime $k$) and with confidence
$1-\delta$. The four architectures (MLP+attention, Neural ODE, Flow Matching,
score-based) are compared on the full metric set above, and the selected proxy
is exported as a differentiable module ($\nabla_\mathbf{F}\widehat{\mathrm{SCR}}$
available) for direct use as an explicit constraint in the Axis~3 programme.

% ----------------------------------------------------------------
\subsection{Axis~3 --- Stochastic Control: Joint Allocation and Deep Hedging}
% ----------------------------------------------------------------

\noindent\textbf{Research question.}
\begin{rqthm}
How to learn a joint stochastic control policy optimising TPA factor exposures
and hedging positions under a differentiable SCR capital constraint, with
regime-adaptive objectives automatically derived from the dominant risk factor?
\end{rqthm}

\noindent\textbf{State of the art.}

\textit{Deep reinforcement learning for ALM.}
Wekwete (2023) demonstrated the applicability of Deep RL (PPO, SAC) to
insurance ALM in stochastic environments, establishing the feasibility of
learning allocation policies in high-dimensional state spaces under regulatory
constraints.

\textit{Deep hedging.}
Buehler et al.\ (2019) introduced deep hedging: learning a hedging strategy
via a neural network minimising a convex risk measure without market
completeness assumptions. The approach has been extended to extreme risk
(Carbonneau, 2021; Cao et al., 2023) and multi-asset portfolios under stress
regimes, and provides the algorithmic backbone for the joint
allocation-hedging objective of Axis~3.

\textit{Stochastic control under dynamic risk constraints.}
R\'{e}veillac (2013) studies CRRA utility maximisation subject to dynamic
risk constraints on trading strategies, characterised by quadratic BSDEs.
Within the class of time-consistent distortion risk measures a three-fund
separation result holds --- directly informing the stochastic control
formulation where the joint policy operates under a dynamic SCR constraint
generated by a general class of risk measures.

\textit{Shortfall systemic risk measures and capital allocation.}
Rosazza Gianin et al.\ (2024) develop a theory of shortfall systemic risk
measures with axiomatic capital allocation rules via a generalised
collapse-to-the-mean framework. Their results connect to the
distortion-based capital constraint of Axis~3 and motivate the use of
$\mathrm{CVaR}_{0.99}$ of the capital deficit as the hedging objective.

\noindent\textbf{Identified gaps.}
No existing framework jointly learns allocation and hedging under: (i)~a
regime-conditional, non-linear dependence model; (ii)~a fast, differentiable
capital constraint; (iii)~private asset engagement restrictions; and (iv)~an
automatically identified portfolio-specific objective function. The full
coupling of Axes~1 and~2 into a stochastic control loop is the contribution
of Axis~3.

\noindent\textbf{Proposed approach.}
A latent objective map $\psi_t = g_\eta(\mathbf{F}_t,z_t,\mathbf{P}_t,
\widehat{\mathrm{SCR}}_t)$ identifies the dominant factor objective per
portfolio and regime. The agent state is
$s_t = (\mathbf{F}_t,\hat{C}_{z_t},z_t,\mathbf{P}_t,\widehat{\mathrm{SCR}}_t)$
and the action $(\mathbf{w}_t,\mathbf{h}_t)$ is the joint
allocation-hedging pair. The policy $\pi$ solves
\begin{equation}
  \max_\pi\;\mathbb{E}_\pi\!\left[\sum_{t=0}^T
  r_{\psi_t}(\mathbf{w}_t,\mathbf{h}_t,\mathbf{F}_{t+1})
  - \mathcal{P}_t \right]
  \quad\text{s.t.}\quad\widehat{\mathrm{SCR}}_t \le \bar{C},
\label{eq:obj}
\end{equation}
where $r_{\psi_t}$ is the regime-and-portfolio-specific objective and
$\mathcal{P}_t$ aggregates hedging cost, turnover and new-illiquid-commitment
penalties. The hedging component targets
$\mathrm{CVaR}_{0.99}$ reduction of the capital deficit, with the
regime-conditional factor dependence of Axis~1 embedded directly in the
simulation environment.


% ================================================================
\section{Scientific Risks and Mitigation}
% ================================================================

The programme is ambitious; its main scientific and operational risks have been
anticipated, and each axis is designed with a fallback that preserves the
overall contribution.

\begin{tabularx}{\linewidth}{@{}p{4.4cm}>{\centering\arraybackslash}p{1cm}>{\raggedright\arraybackslash}X@{}}
\toprule
\textbf{Risk} & \textbf{Lik.} & \textbf{Mitigation / fallback} \\
\midrule
cGAN mode collapse or unstable training (Axis~1)
& Med.
& The CVAE and conditional SGM are trained in parallel as substitutes;
  the SGM enjoys formal $W_2$ guarantees and is the primary fallback if
  adversarial training fails. \\
\addlinespace
Conditional $W_2$ bound not provable in full generality (Axis~1)
& Med.
& A partial result under piecewise log-concavity (Gaussian-mixture regimes,
  following Gentiloni-Silveri \& Ocello, 2025) already covers the practically
  relevant multi-regime case. \\
\addlinespace
Limited access to granular BEL / private-asset data
& Med.
& The actuarial design-of-experiment is augmented synthetically via TabDDPM;
  the proxy is validated on historical crises (2008, 2020, 2022). \\
\addlinespace
Proxy accuracy degrades in extreme stress zones (Axis~2)
& Low
& Adaptive sample weighting concentrates training on decision-relevant tails;
  MLMC control variates (Bourgey, 2020) bound the residual $L^2$ error. \\
\addlinespace
RL policy instability / non-convergence (Axis~3)
& Med.
& The differentiable proxy enables direct gradient-based control as a
  fallback to model-free RL; deep hedging provides a validated warm start. \\
\bottomrule
\end{tabularx}

% ================================================================
\section{Timeline (36 Months)}
% ================================================================

\begin{itemize}[leftmargin=1.6em,itemsep=2pt]
  \item \textbf{Year~1 --- Axis~1.}
        Bloomberg factor database with NAV lags; online HMM-SMC calibration;
        regime-conditional GAN/CVAE validated on 2008/2020/2022 crises.
        \emph{Submission: Journal of Multivariate Analysis / Finance \& Stochastics.}
  \item \textbf{Year~2 --- Axis~2.}
        Actuarial design-of-experiment (Prophet/MoSes) with TabDDPM
        augmentation; SCR proxy architecture comparison and Volterra $L^2$
        bounds; differentiable proxy embedded as optimisation constraint.
        \emph{Submission: Insurance: Mathematics and Economics.}
  \item \textbf{Year~3 --- Axis~3.}
        Joint allocation-and-hedging agent with automatic objective
        identification; full three-axis coupling; multi-regime out-of-sample
        backtesting; defence. \emph{Submission: Quantitative Finance / JFQA.}
\end{itemize}

% ================================================================
\section{Conclusion}
% ================================================================

This thesis delivers a tightly coupled three-module system designed to make the
\textbf{Total Portfolio Approach operational for a life insurer under
Solvency~II}, with direct integration into Generali France's proprietary
platform GPS:

\begin{itemize}[leftmargin=1.6em,itemsep=3pt]
  \item \textbf{Module 1 (Axis~1)} replaces static ESG correlations with an
        auditable regime-conditional generative model, enabling real-time
        scenario generation that reflects the actual dependence structure of
        the prevailing macroeconomic regime, including the path-dependent
        correction dynamics of private assets.

  \item \textbf{Module 2 (Axis~2)} replaces the multi-hour SCR calculation
        with a $<5$ms differentiable proxy grounded in Volterra functional
        theory, making dynamic portfolio optimisation loops computationally
        feasible and the capital constraint analytically tractable as an
        explicit gradient-available term.

  \item \textbf{Module 3 (Axis~3)} replaces the annual, siloed allocation
        process with a dynamic joint policy that simultaneously optimises TPA
        factor exposures and hedging positions under an automatically
        identified portfolio-specific objective and an explicit SCR constraint.
\end{itemize}

Beyond the industrial deliverables, the thesis makes three mathematically
self-contained contributions: the extension of $W_2$ convergence theory to
regime-conditional score-based models under sequential estimation error; the
formalisation of private asset valuation bias through the theory of Volterra
processes and path-dependent risk measures; and the characterisation of
optimal stochastic control policies under dynamic, non-Markovian capital
constraints derived from a general class of shortfall risk measures. Together,
these contributions establish rigorous foundations for a new generation of
quantitative tools in insurance balance sheet management --- tools that are
simultaneously \emph{theoretically grounded}, \emph{computationally efficient},
\emph{interoperable} and \emph{regulatory-compliant}.

\medskip
\noindent\textbf{Beyond the three years.}
The certified-generation / differentiable-proxy / stochastic-control triptych
is deliberately generic. The same architecture extends naturally to dynamic
ORSA projection, climate stress-testing under regime-dependent transition
scenarios, and IFRS~17 risk-adjustment computation --- opening a longer-term
research programme on \emph{certified, capital-aware decision systems} for the
insurance balance sheet.

% ================================================================
\section*{References}
% ================================================================
\renewcommand{\section}[2]{}

\begin{thebibliography}{30}

\bibitem{abijabar2019}
Abi Jaber, E., Larsson, M., \& Pulido, S. (2019).
Affine Volterra processes.
\emph{The Annals of Applied Probability}, 29(6), 3155--3200.

\bibitem{bardoscia2021}
Bardoscia, M., Barucca, P., Battiston, S., Caccioli, F., Cimini, G.,
Garlaschelli, D., Saracco, F., Squartini, T., \& Caldarelli, G. (2021).
The physics of financial networks.
\emph{Nature Reviews Physics}, 3(7), 490--507.

\bibitem{bourgey2020}
Bourgey, F. (2020).
\emph{Stochastic approximations for financial risk computations}.
PhD thesis.

\bibitem{buehler2019}
Buehler, H., Gonon, L., Teichmann, J., \& Wood, B. (2019).
Deep hedging.
\emph{Quantitative Finance}, 19(8), 1271--1291.

\bibitem{caia2024}
CAIA Association. (2024).
\emph{Innovation Unleashed: The Rise of Total Portfolio Approach}.

\bibitem{cao2023}
Cao, J., Chen, J., Hull, J., \& Poulos, Z. (2023).
Deep hedging of derivatives using reinforcement learning.
\emph{The Journal of Financial Data Science}, 5(1), 10--27.

\bibitem{carbonneau2021}
Carbonneau, A. (2021).
Deep hedging of long-term financial derivatives.
\emph{Insurance: Mathematics and Economics}, 99, 327--340.

\bibitem{chen2018}
Chen, R.T.Q., Rubanova, Y., Bettencourt, J., \& Duvenaud, D. (2018).
Neural ordinary differential equations.
\emph{Advances in Neural Information Processing Systems}, 31.

\bibitem{chopin2020}
Chopin, N., \& Papaspiliopoulos, O. (2020).
\emph{An Introduction to Sequential Monte Carlo}. Springer.

\bibitem{demarco2025}
De Marco, S., Pham, H., \& Zanni, D. (2025).
Schr\"{o}dinger bridges with jumps for time series generation.
\emph{arXiv preprint arXiv:2602.20011}.

\bibitem{eleuch2019}
El Euch, O., \& Rosenbaum, M. (2019).
The characteristic function of rough Heston models.
\emph{Mathematical Finance}, 29(1), 3--38.

\bibitem{fermanian2013}
Fermanian, J.-D. (2013).
An overview of the goodness-of-fit test problem for copulas.
In \emph{Copulae in Mathematical and Quantitative Finance},
Springer, 61--89.

\bibitem{fermanian2022}
Fermanian, J.-D. (2022).
Conditional copulas, dependence modeling and the role of conditioning.
\emph{Annals of Statistics}.

\bibitem{ocello2025icml}
Gentiloni-Silveri, M., \& Ocello, A. (2025).
Beyond log-concavity and score regularity: Improved convergence bounds
for score-based generative models in $W_2$-distance.
\emph{Forty-Second International Conference on Machine Learning (ICML)}.

\bibitem{goodfellow2014}
Goodfellow, I. et al. (2014).
Generative adversarial networks.
\emph{Advances in Neural Information Processing Systems}, 27.

\bibitem{hamilton1989}
Hamilton, J.D. (1989).
A new approach to the economic analysis of nonstationary time series.
\emph{Econometrica}, 57(2), 357--384.

\bibitem{huber2022}
Huber, M., \& Schmocker, B. (2022).
A neural network approach to efficient valuation of large portfolios
of variable annuities.
\emph{Insurance: Mathematics and Economics}, 105, 96--116.

\bibitem{kotelnikov2023}
Kotelnikov, A., Baranchuk, D., Rubachev, I., \& Babenko, A. (2023).
TabDDPM: Modelling Tabular Data with Diffusion Models.
\emph{International Conference on Machine Learning (ICML)}.

\bibitem{krah2018}
Krah, A.-S., Nikolic, Z., \& Korn, R. (2018).
A Least-Squares Monte Carlo framework in proxy modeling of life
insurance companies.
\emph{Risks}, 6(2), 62.

\bibitem{lipman2022}
Lipman, Y., Chen, R.T.Q., Ben-Hamu, H., Nickel, M., \& Le, M. (2022).
Flow Matching for Generative Modeling.
\emph{International Conference on Learning Representations (ICLR)}.

\bibitem{nystrup2017}
Nystrup, P., Madsen, H., \& Lindstrom, E. (2017).
Dynamic portfolio optimization across hidden market regimes.
\emph{Quantitative Finance}, 18(1), 83--95.

\bibitem{pham2009}
Pham, H. (2009).
\emph{Continuous-Time Stochastic Control and Optimization
with Financial Applications}. Springer.

\bibitem{reveillac2013}
R\'{e}veillac, A. (2013).
CRRA utility maximization under dynamic risk constraints.
\emph{Mathematical Finance}.

\bibitem{rosazza2024}
Rosazza Gianin, E., et al.\ (2024).
Shortfall systemic risk measures, capital allocation rules and
generalised collapse to the mean: Theory and practice.
\emph{Working paper}.

\bibitem{sklar1959}
Sklar, A. (1959).
Fonctions de repartition a $n$ dimensions et leurs marges.
\emph{Publications de l'Institut de Statistique de Paris}, 8, 229--231.

\bibitem{ocello2024tmlr}
Strasman, S., Ocello, A., Boyer, C., Le~Corff, S., \& Lemaire, V. (2024).
An analysis of the noise schedule for score-based generative models.
\emph{Transactions on Machine Learning Research (TMLR)}.

\bibitem{thinking_ahead}
Thinking Ahead Institute. (2025).
\emph{TPA Report: From Vision to Execution}. WTW / CAIA.

\bibitem{wekwete2023}
Wekwete, T. (2023).
Application of Deep Reinforcement Learning in Asset Liability Management.
\emph{Actuarial Society of South Africa Convention Paper}.

\bibitem{wiese2020}
Wiese, M., Knobloch, R., Korn, R., \& Kretschmer, P. (2020).
Quant GANs: Deep generation of financial time series.
\emph{Quantitative Finance}, 20(9), 1419--1440.

\bibitem{yoon2019}
Yoon, J., Jarrett, D., \& van der Schaar, M. (2019).
Time-series generative adversarial networks.
\emph{Advances in Neural Information Processing Systems}, 32.

\end{thebibliography}

\end{document}
